\section{Reproducibility and Benchmark Hygiene}
\label{app:hygiene}

\draft{Prompt-set manifests (SHA-256 over problem IDs) for every eval and
training set; avg@$k$ seeds; sandbox image hash for the code reward.
HELMET note from the status report, kept for \cref{app:helmet}: Our HELMET port initially drew the dev/eval
split unseeded (dropped the harness-level seeding HELMET itself uses); found
by cross-host replication, fixed in PR \#145. Cross-host splits shared only
8--10/100 questions, but all campaign numbers came from one host and one
cache, so paired comparisons stand. Canonical split caches (8k campaign draw
and seeded 32k draws) are published with SHA-256 manifests over record IDs.
Upstream HELMET has a smaller demo-sampling nondeterminism; issue filed
(princeton-nlp/HELMET\#43) with a fix.}

\section{Additional Memory Measurements}
\label{app:memory}

\draft{Qwen2.5-3B-Instruct on A10G 24\,GB (from
\texttt{docs/GRPO\_CONTEXT\_SCALING.md}): dense 14.3 / 16.4 / 19.8\,GB at 2k /
8k / 16k, OOM at 24k and 32k; H2O ($K=4096$) flat at 14.9\,GB from 2k to 32k
(re-measured in a fresh process). Note the in-process contamination caveat
and that per-config peaks should be measured in fresh processes.}

\section{Training Details}
\label{app:training}

\draft{Full hyperparameters per run: learning rate, group size, steps,
gradient accumulation, chunk size, layers per cell, optimiser (8-bit paged
AdamW), reward definition, prompt templates. Point to the exact config files
and commit hashes in the repository.}

\section{Reward Definitions}
\label{app:rewards}

\draft{Math: final-answer extraction and equivalence check (which verifier:
math-verify / Polaris's). Code: sandbox, time limits, pass-all-tests binary
reward, following Liu and Zhang (2025) as Sparrow does.}

\section{Long-Input QA Results (HELMET)}
\label{app:helmet}

\draft{Supporting evidence from the earlier campaign: 0.5B@8k GRPO/RLOO/SFT
dense vs evicting on NaturalQuestions and HotpotQA (parity within noise, RL
$>$ SFT, RLOO+EM nq 0.43 vs 0.33 base); 7B@32k base probes and RLOO runs;
the \texttt{em\_f1} reward and the style-drift artifact. Different task
family from the main text, same conclusion on parity.}

\section{Comparison of Experimental Protocols with Sparrow}
\label{app:sparrow}

\draft{Keep this table for the rebuttal even if it does not make the main
text.}

\begin{table}[h]
  \caption{Protocol comparison with Sparrow \citep{sadhukhan2026sparrow}.}
  \label{tab:sparrow_protocol}
  \centering\small
  \begin{tabular}{p{0.2\linewidth}p{0.38\linewidth}p{0.34\linewidth}}
    \toprule
    & Sparrow & This work \\
    \midrule
    Where sparsity acts & rollout only (dense policy) & rollout and gradient (sparse policy) \\
    Sparsity mechanism & block-sparse page selection, dynamic schedule & H2O eviction + int8 KV, fixed budget \\
    Generation cutoff & 24--37k & \draft{TODO} \\
    Models & Qwen3 thinking 1.7B/4B/8B/14B & Qwen3 thinking 1.7B/4B (\draft{8B}) \\
    Train data & Polaris (math), TACO 21K (code), DeepScaleR (warmup) & same \\
    Eval & AIME 24/25/26, AMC 23/24, MATH500; LiveCodeBench, HumanEval+, MBPP+ & same \\
    Hardware & 4$\times$H200 up to 32$\times$H200 & 1$\times$L40S 48\,GB (fleet of single-GPU workers) \\
    Headline metric & rollout speedup at dense-parity accuracy & memory frontier; accuracy where dense OOMs \\
    Stability control & tail acceptance-rate threshold ($\tau\approx0.86$) & none needed (sampler $=$ policy) \\
    \bottomrule
  \end{tabular}
\end{table}
